\begin{table}[!htbp]
\centering
\caption{Baseline characteristics of high-premium-field value-added compliers}
\label{tab:high-premium-field-va-complier-ses}
\begin{threeparttable}
\small
\setlength{\tabcolsep}{4pt}
\begin{tabular}{lcccc}
\toprule
 & SAE mean & Above P25 & Above median & Above P75 \\
\midrule
Baseline math score & -0.233 & +0.053** & +0.088*** & +0.084*** \\
 &  & (0.024) & (0.017) & (0.015) \\
Baseline verbal score & -0.189 & +0.071*** & +0.083*** & +0.062*** \\
 &  & (0.024) & (0.017) & (0.015) \\
Income decile & 4.519 & +0.121* & +0.227*** & +0.314*** \\
 &  & (0.063) & (0.045) & (0.040) \\
Father's education (years) & 10.771 & +0.046 & +0.188*** & +0.321*** \\
 &  & (0.071) & (0.050) & (0.045) \\
Mother's education (years) & 10.941 & +0.152** & +0.251*** & +0.330*** \\
 &  & (0.067) & (0.048) & (0.043) \\
\bottomrule
\end{tabular}
\begin{tablenotes}[flushleft]
\footnotesize
\item Notes: Each threshold column reports the estimated mean baseline characteristic among students induced by their first-round offer to attend a school above the stated high-premium-field value-added threshold, minus the corresponding mean among all timely SAE applicants in the 2018--2020 cohorts. The SAE mean is common across thresholds. Parentheses contain heteroskedasticity-robust standard errors that account for estimation of the complier and population means and their covariance. High value added is defined using equally weighted school quantiles of EB-shrunken value added. The complier population and estimation sample may differ across thresholds. Differences describe baseline composition and are not treatment effects. $^{*}p<0.10$, $^{**}p<0.05$, $^{***}p<0.01$.
\end{tablenotes}
\end{threeparttable}
\end{table}
